\documentclass[10pt,a4paper]{amsart}
\usepackage{amsmath,amssymb,amsthm,graphicx}
% Retained preamble for the Stacks Project chapter snapshot in this directory.
% The chapter's own preamble.tex is a copy of the shared Stacks Project preamble; the block
% below is sliced MECHANICALLY, line for line, from it. Nothing was added, removed or
% redefined; the wrapper class, the xr-hyper/hyperref wiring and the externaldocument list
% are omitted because the assembler supplies the document class.
%
% lines taken from preamble.tex: 3, 5, 9, 11, 12, 156, 157, 158, ... (51 lines total)
%
% package coverage restored from the chapter's own preamble (2026-10-13): xy, xy-2cell

\usepackage[all]{xy}\usepackage{verbatim}
\xyoption{2cell}
\UseAllTwocells

\newenvironment{reference}{\comment}{\endcomment}
\newenvironment{slogan}{\comment}{\endcomment}
\newenvironment{history}{\comment}{\endcomment}
\theoremstyle{plain}
\newtheorem{theorem}[subsection]{Theorem}
\newtheorem{proposition}[subsection]{Proposition}
\newtheorem{lemma}[subsection]{Lemma}
\theoremstyle{definition}
\newtheorem{definition}[subsection]{Definition}
\newtheorem{example}[subsection]{Example}
\newtheorem{exercise}[subsection]{Exercise}
\newtheorem{situation}[subsection]{Situation}
\theoremstyle{remark}
\newtheorem{remark}[subsection]{Remark}
\newtheorem{remarks}[subsection]{Remarks}
\numberwithin{equation}{subsection}
\def\lim{\mathop{\mathrm{lim}}\nolimits}
\def\colim{\mathop{\mathrm{colim}}\nolimits}
\def\Spec{\mathop{\mathrm{Spec}}}
\def\Hom{\mathop{\mathrm{Hom}}\nolimits}
\def\Ext{\mathop{\mathrm{Ext}}\nolimits}
\def\SheafHom{\mathop{\mathcal{H}\!\mathit{om}}\nolimits}
\def\SheafExt{\mathop{\mathcal{E}\!\mathit{xt}}\nolimits}
\def\Sch{\mathit{Sch}}
\def\Mor{\mathop{\mathrm{Mor}}\nolimits}
\def\Ob{\mathop{\mathrm{Ob}}\nolimits}
\def\Sh{\mathop{\mathit{Sh}}\nolimits}
\def\NL{\mathop{N\!L}\nolimits}
\def\CH{\mathop{\mathrm{CH}}\nolimits}
\def\proetale{{pro\text{-}\acute{e}tale}}
\def\etale{{\acute{e}tale}}
\def\QCoh{\mathit{QCoh}}
\def\Ker{\mathop{\mathrm{Ker}}}
\def\Im{\mathop{\mathrm{Im}}}
\def\Coker{\mathop{\mathrm{Coker}}}
\def\Coim{\mathop{\mathrm{Coim}}}
\DeclareMathSymbol{\boxtimes}{\mathbin}{AMSa}{"02}
\def\QCohstack{\mathcal{QC}\!\mathit{oh}}
\def\Cohstack{\mathcal{C}\!\mathit{oh}}
\def\Spacesstack{\mathcal{S}\!\mathit{paces}}
\def\Quotfunctor{\mathrm{Quot}}
\def\Hilbfunctor{\mathrm{Hilb}}
\def\Curvesstack{\mathcal{C}\!\mathit{urves}}
\def\Polarizedstack{\mathcal{P}\!\mathit{olarized}}
\def\Complexesstack{\mathcal{C}\!\mathit{omplexes}}
\def\Pic{\mathop{\mathrm{Pic}}\nolimits}
\def\Picardstack{\mathcal{P}\!\mathit{ic}}
\def\Picardfunctor{\mathrm{Pic}}
\def\Deformationcategory{\mathcal{D}\!\mathit{ef}}
\usepackage{cleveref}
\crefname{originallemma}{original lemma}{original lemmas}
\Crefname{originallemma}{Original lemma}{Original lemmas}
\makeatletter
\newlabel{groupoids-lemma-etale}{{0DT9}{source}{}{sourcepointer-1}{}}
\newlabel{groupoids-lemma-etale@cref}{{[originallemma][0][]0DT9}{[1][0][]source}{}{}{}}
\makeatother
\crefname{originallemma}{original lemma}{original lemmas}
\Crefname{originallemma}{Original lemma}{Original lemmas}
\makeatletter
\newlabel{groupoids-quotients-lemma-base-change-quotient-stack}{{0DTG}{source}{}{sourcepointer-2}{}}
\newlabel{groupoids-quotients-lemma-base-change-quotient-stack@cref}{{[originallemma][0][]0DTG}{[1][0][]source}{}{}{}}
\makeatother
\crefname{originallemma}{original lemma}{original lemmas}
\Crefname{originallemma}{Original lemma}{Original lemmas}
\makeatletter
\newlabel{lemma-affine-over-well-nigh-affine}{{lemma-affine-over-well-nigh-affine}{source}{}{sourcepointer-3}{}}
\newlabel{lemma-affine-over-well-nigh-affine@cref}{{[originallemma][0][]lemma-affine-over-well-nigh-affine}{[1][0][]source}{}{}{}}
\makeatother
\crefname{originallemma}{original lemma}{original lemmas}
\Crefname{originallemma}{Original lemma}{Original lemmas}
\makeatletter
\newlabel{lemma-well-nigh-affine}{{lemma-well-nigh-affine}{source}{}{sourcepointer-4}{}}
\newlabel{lemma-well-nigh-affine@cref}{{[originallemma][0][]lemma-well-nigh-affine}{[1][0][]source}{}{}{}}
\makeatother
\crefname{originallemma}{original lemma}{original lemmas}
\Crefname{originallemma}{Original lemma}{Original lemmas}
\makeatletter
\newlabel{lemma-well-nigh-affine-moduli-space}{{lemma-well-nigh-affine-moduli-space}{source}{}{sourcepointer-5}{}}
\newlabel{lemma-well-nigh-affine-moduli-space@cref}{{[originallemma][0][]lemma-well-nigh-affine-moduli-space}{[1][0][]source}{}{}{}}
\makeatother
\crefname{originallemma}{original lemma}{original lemmas}
\Crefname{originallemma}{Original lemma}{Original lemmas}
\makeatletter
\newlabel{stacks-morphisms-lemma-aut-iso-unramified}{{0DU9}{source}{}{sourcepointer-6}{}}
\newlabel{stacks-morphisms-lemma-aut-iso-unramified@cref}{{[originallemma][0][]0DU9}{[1][0][]source}{}{}{}}
\makeatother
\crefname{originallemma}{original lemma}{original lemmas}
\Crefname{originallemma}{Original lemma}{Original lemmas}
\makeatletter
\newlabel{stacks-morphisms-lemma-stabilizer-preserving}{{0DU7}{source}{}{sourcepointer-7}{}}
\newlabel{stacks-morphisms-lemma-stabilizer-preserving@cref}{{[originallemma][0][]0DU7}{[1][0][]source}{}{}{}}
\makeatother
\crefname{originallemma}{original lemma}{original lemmas}
\Crefname{originallemma}{Original lemma}{Original lemmas}
\makeatletter
\newlabel{stacks-morphisms-remark-identify-automorphism-groups}{{0DTW}{source}{}{sourcepointer-8}{}}
\newlabel{stacks-morphisms-remark-identify-automorphism-groups@cref}{{[originallemma][0][]0DTW}{[1][0][]source}{}{}{}}
\makeatother
\title{Etale morphisms of well-nigh-affine stacks base-change the invariant-ring moduli space}
\author{The Stacks Project Authors}
\date{stacks-project:stacks-more-morphisms.tex:local-snapshot-20261012; GNU FDL 1.2 or later}
\begin{document}
\maketitle
\noindent\textit{Editorial note.} Statement span is lines 3516-3533 of sources/stacks-stacks-more-morphisms/pinned-original-source.tex; proof span is lines 3535-3589, the complete author proof (53 body lines).\par\medskip
\noindent\textit{Editorial note.} Navigation stubs: 8 cited-but-not-included labels are supplied as non-load-bearing pointers with their pinned tags where the snapshot tag index has them.\par\medskip
\noindent\textit{Editorial note.} extra\_packages is empty on purpose: the retained preamble already loads xy with options and a second declaration clashes in the assembler wrapper.\par\medskip
\noindent\textit{Editorial note.} Text is retained verbatim, including the author's own typos and placeholders.\par\medskip
\noindent\textit{Editorial note.} Queue index 25 of queue\_stacks-stacks-more-morphisms.json (pb=2059, stmt 3516-3533, proof 3535-3589).\par\medskip
\noindent\textit{Editorial note.} stacks-stacks-more-morphisms wave smm01 (2026-10-14): fresh adjudication accept per lane-s3/smm\_judgments.md; pool per Chief 10:5x open-pool ruling (net 59 with-proof rows, bloodline lane-s3 by anchor audit, v2 true-key namespace stacks:stacks-more-morphisms.tex, row lemma-deform-property-fp-over-ft avoided on cross-chapter key collision); 60 with-proof rows judged 45 accept / 14 reject / 1 avoid; patch-sd-envs contract spans scanned.\par\medskip
\noindent\textit{Allowed standard prerequisites (editorial).}\begin{itemize}
\item Explicit prerequisite boundary. The cited inputs are prerequisites; the argument in the retained span is not replaced by any cited result.
\end{itemize}\medskip
\ifdefined\hypertarget\hypertarget{sourcepointer-1}{}\fi
\noindent\textit{Source reference.} Original lemma 0DT9: sources/stacks-stacks-more-morphisms/pinned-original-source.tex, source lines 3589--3589. Cited by the retained statement or proof but not part of the collected spans; supplied only as a navigation pointer. The argument does not depend on its content.\par\medskip
\ifdefined\hypertarget\hypertarget{sourcepointer-2}{}\fi
\noindent\textit{Source reference.} Original lemma 0DTG: sources/stacks-stacks-more-morphisms/pinned-original-source.tex, source lines 3589--3589. Cited by the retained statement or proof but not part of the collected spans; supplied only as a navigation pointer. The argument does not depend on its content.\par\medskip
\ifdefined\hypertarget\hypertarget{sourcepointer-3}{}\fi
\noindent\textit{Source reference.} Original lemma lemma-affine-over-well-nigh-affine: sources/stacks-stacks-more-morphisms/pinned-original-source.tex, source lines 3396--3402. Cited by the retained statement or proof but not part of the collected spans; supplied only as a navigation pointer. The argument does not depend on its content.\par\medskip
\ifdefined\hypertarget\hypertarget{sourcepointer-4}{}\fi
\noindent\textit{Source reference.} Original lemma lemma-well-nigh-affine: sources/stacks-stacks-more-morphisms/pinned-original-source.tex, source lines 3339--3346. Cited by the retained statement or proof but not part of the collected spans; supplied only as a navigation pointer. The argument does not depend on its content.\par\medskip
\ifdefined\hypertarget\hypertarget{sourcepointer-5}{}\fi
\noindent\textit{Source reference.} Original lemma lemma-well-nigh-affine-moduli-space: sources/stacks-stacks-more-morphisms/pinned-original-source.tex, source lines 3425--3433. Cited by the retained statement or proof but not part of the collected spans; supplied only as a navigation pointer. The argument does not depend on its content.\par\medskip
\ifdefined\hypertarget\hypertarget{sourcepointer-6}{}\fi
\noindent\textit{Source reference.} Original lemma 0DU9: sources/stacks-stacks-more-morphisms/pinned-original-source.tex, source lines 3589--3589. Cited by the retained statement or proof but not part of the collected spans; supplied only as a navigation pointer. The argument does not depend on its content.\par\medskip
\ifdefined\hypertarget\hypertarget{sourcepointer-7}{}\fi
\noindent\textit{Source reference.} Original lemma 0DU7: sources/stacks-stacks-more-morphisms/pinned-original-source.tex, source lines 3589--3589. Cited by the retained statement or proof but not part of the collected spans; supplied only as a navigation pointer. The argument does not depend on its content.\par\medskip
\ifdefined\hypertarget\hypertarget{sourcepointer-8}{}\fi
\noindent\textit{Source reference.} Original lemma 0DTW: sources/stacks-stacks-more-morphisms/pinned-original-source.tex, source lines 3589--3589. Cited by the retained statement or proof but not part of the collected spans; supplied only as a navigation pointer. The argument does not depend on its content.\par\medskip
\begin{lemma}
\label{lemma-well-nigh-affine-moduli-space-etale}
Let $h : \mathcal{X}' \to \mathcal{X}$ be a morphism of algebraic stacks.
Assume $\mathcal{X}'$ and $\mathcal{X}$ are well-nigh affine,
$h$ is \'etale, and $h$ induces isomorphisms on automorphism groups
(Morphisms of Stacks, Remark
\ref{stacks-morphisms-remark-identify-automorphism-groups}).
Then there exists a cartesian diagram
$$
\xymatrix{
\mathcal{X}' \ar[d] \ar[r] & \mathcal{X} \ar[d] \\
M' \ar[r] & M
}
$$
where $M' \to M$ is \'etale and
the vertical arrows are the moduli spaces constructed in
Lemma \ref{lemma-well-nigh-affine-moduli-space}.
\end{lemma}

\begin{proof}
Observe that $h$ is representable by algebraic spaces by
Morphisms of Stacks, Lemmas
\ref{stacks-morphisms-lemma-aut-iso-unramified} and
\ref{stacks-morphisms-lemma-stabilizer-preserving}.
Choose an affine scheme $U$ and a
surjective, flat, finite, and finitely presented morphism $U \to \mathcal{X}$.
Then $U' = \mathcal{X}' \times_\mathcal{X} U$ is an algebraic
space with a finite (in particular affine) morphism $U' \to \mathcal{X}'$.
By Lemma \ref{lemma-affine-over-well-nigh-affine}
we conclude that $U'$ is affine.
Setting $R = U \times_\mathcal{X} U$ and $R' = U' \times_{\mathcal{X}'} U'$
we obtain groupoids $(U, R, s, t, c)$ and $(U', R', s', t', c')$
such that $\mathcal{X} = [U/R]$ and $\mathcal{X}' = [U'/R']$,
see proof of Lemma \ref{lemma-well-nigh-affine}.
we see that the diagrams
$$
\xymatrix{
R' \ar[d]_{s'} \ar[r]_f & R \ar[d]^s \\
U' \ar[r]^f & U
}
\quad
\quad
\xymatrix{
R' \ar[d]_{t'} \ar[r]_f & R \ar[d]^t \\
U' \ar[r]^f & U
}
\quad
\quad
\xymatrix{
G' \ar[d] \ar[r]_f & G \ar[d] \\
U' \ar[r]^f & U
}
$$
are cartesian where $G$ and $G'$ are the stabilizer group schemes.
This follows for the first two by transitivity of fibre products
and for the last one this follows because it is the pullback of the
isomorphism $\mathcal{I}_{\mathcal{X}'} \to
\mathcal{X}' \times_\mathcal{X} \mathcal{I}_\mathcal{X}$
(by the already used Morphisms of Stacks, Lemma
\ref{stacks-morphisms-lemma-aut-iso-unramified}).
Recall that $M$, resp.\ $M'$ was constructed in
Lemma \ref{lemma-well-nigh-affine-moduli-space}
as the spectrum of the ring of $R$-invariant functions on $U$,
resp.\ the ring of $R'$-invariant functions on $U'$.
Thus $M' \to M$ is \'etale and $U' = M' \times_M U$
by Groupoids, Lemma \ref{groupoids-lemma-etale}.
It follows that $R' = M' \times_M U$, in other words
the groupoid $(U', R', s', t', c')$ is the base change of
$(U, R, s, t, c)$ by $M' \to M$.
This implies that the diagram in the lemma is
cartesian by
Quotients of Groupoids, Lemma
\ref{groupoids-quotients-lemma-base-change-quotient-stack}.
\end{proof}

\end{document}
